\documentclass{article}
\usepackage{iclr2027_conference,times}
\input{math_commands.tex}
\usepackage{hyperref}
\usepackage{url}
\usepackage{booktabs}
\usepackage{amsmath}
\usepackage{xcolor}

% ── Draft mode ───────────────────────────────────────────────────────────────
% \draftmodetrue  → internal draft: \draftnote{...} renders in red, unfinished results are visibly marked.
% \draftmodefalse → submission build: every \draftnote is an ERROR (scripts/check_submission.py also fails), so no
%                    placeholder can survive into the PDF. Never ship a PDF built with \draftmodetrue.
\newif\ifdraftmode
\draftmodetrue
\newcommand{\draftnote}[1]{\ifdraftmode{\color{red}\textbf{[DRAFT: #1]}}\else\PackageError{main}{draftnote left in submission build: #1}{}\fi}
% Anonymous submission: keep \iclrfinalcopy commented out. Author block below is the template placeholder on purpose.

\title{Does post-execution feedback improve the choice between stopping, self-revision, expert review and regeneration? \\
A pre-registered, cost-accounted study on BigCodeBench}

\author{Anonymous authors \\ Paper under double-blind review}
% \iclrfinalcopy

\begin{document}
\maketitle

\begin{abstract}
\draftnote{Abstract is written after the pre-registered analysis (decision package, 2026-09-17). Do not register a placeholder. The
registered abstract must contain only results verified in \texttt{evidence/week2} at registration time; see
\texttt{evidence/iclr\_2027/decision\_20260917.md} for the two admissible variants.}
\end{abstract}

\section{Introduction}
\label{sec:intro}
A code-generating agent that has just produced a candidate and run the visible checks faces a decision: submit the candidate as is
(\emph{stop}), ask the same model to revise it (\emph{self-revise}), hand it to a stronger model for review (\emph{expert review}), or
discard the state and let the stronger model \emph{regenerate} from scratch. Each continuation costs money, and the information that
could justify it---the executed check outcome, the failure count, the error class, the size of the artifact---is only available
\emph{after} paying for the first execution.

We ask one question:
\begin{quote}
\emph{Does post-execution observable feedback improve the choice among stop / self-revise / review-by-another-model / regenerate on
\textbf{new} tasks, and does the gain survive once the cost of acquiring that state and the cost of the executions used to learn the
policy are counted?}
\end{quote}
We answer it with a two-stage, pre-registered design on BigCodeBench~\citep{zhuo2024bigcodebench} (Section~\ref{sec:design}):
an \emph{exploration} set on which every action is executed and a small offline policy is learned and frozen, and a disjoint
\emph{confirmation} set on which the frozen policy is compared, task by task, against the single best fixed action chosen on the
exploration set. The primary estimand and the decision rule were fixed before any confirmation task was executed.

\paragraph{Scope of the claims.} The policy is an offline routing rule fitted on exploration outcomes; the LLM weights are frozen
throughout, but the router \emph{is} learned from labelled executions, and we count those executions as a cost. We do not claim
training-free operation, online self-improvement, novelty of the action set, or superiority over routers we did not run
(Section~\ref{sec:related}). Conclusions are restricted to the eligible BigCodeBench pool and the two models used.
\draftnote{Numbers and the contribution list are filled from the decision package only.}

\section{Problem definition}
\label{sec:problem}
Let a task $i$ have a first-pass candidate $s_0(i)$ produced by model $A$, together with the executed visible-check result. A
\emph{continuation arm} $a\in\mathcal{A}=\{\text{T},\text{A-self},\text{A-self}{\times}k,\text{A-role},\text{B-expert},\text{B-solo}\}$
maps the state to a final artifact: T submits $s_0$; A-self and A-self${\times}k$ ask $A$ to revise once or up to $k$ times
(stopping on visible pass); A-role asks $A$ in a reviewer role; B-expert asks the stronger model $B$ to review and correct; B-solo
discards $s_0$ and regenerates with $B$. The hidden outcome $y_{i,a}\in\{0,1\}$ is the official hidden test result of the final
artifact, obtained post hoc and never shown to any model.

A \emph{state-conditional policy} $\pi(x_i)\in\mathcal{A}$ chooses an arm from decision-time features $x_i$. We distinguish
\emph{pre} features (computable from the task prompt alone) from \emph{post} features (computable only after $s_0$ was generated and
checked). The quantity of interest is the \emph{policy gain over the best fixed strategy},
\begin{equation}
\Delta(\pi,\hat a)=\mathbb{E}_i\big[y_{i,\pi(x_i)}-y_{i,\hat a}\big],
\end{equation}
where $\hat a$ is the fixed arm with the highest exploration success rate. Three quantities must not be confused: (1) the
\emph{true-probability oracle gain} $\mathbb{E}_i[\max_a q_{ia}]-\max_a\mathbb{E}_i[q_{ia}]$, an upper bound on any policy gain that is
observable only in simulation; (2) the \emph{in-sample maximum of single realisations} $\frac1n\sum_i\max_a y_{ia}-\max_a\frac1n\sum_i y_{ia}$,
which contains realisation luck and is reported only as a descriptive statistic; (3) the gain of a policy learned from observable
features, which is what we test.

\section{Method: an offline, frozen routing policy}
\label{sec:method}
\paragraph{Features.} Pre features: prompt length, number of docstring examples, argument count, presence of a Returns clause, and
one-hot library requirements (12 most frequent in the eligible pool + other). Post features: visible outcome, number of failed checks,
error class (8 most frequent + other/none), artifact size, AST statistics (functions, try-blocks, imports), and the token/dollar cost of
$s_0$. Hidden results are never features. \draftnote{cite features.py version FEATURES\_VERSION=features-i1; list is fixed.}

\paragraph{Learners.} P2 fits one $L_2$-regularised logistic model per arm, $\hat p_a(x)$, and picks $\arg\max_a\hat p_a(x)$; P2-post
uses pre+post features, P2-pre uses pre features only, each with its own feature subset and $\lambda$ selected on its own menu. P1 is a
lookup table on the visible-outcome category and serves as a diagnostic. Feature subsets (at most 20 of 28 columns) and
$\lambda\in\{0.3,1,3\}$ are chosen by nested cross-validation inside the exploration set (selection inside each training fold); the
frozen model is refit once on the full exploration set. \draftnote{The exact procedure is \texttt{gap\_analysis.explore\_fit}; the same
function drives the simulator (Appendix).}

\paragraph{Freezing.} After exploration, the feature subsets, $\lambda$, coefficients, the P1 table, $\hat a$, the continuation budget
$b_{\mathrm{cont}}$ and the seed are written to a single file whose hash is recorded before any confirmation task is run.

\section{Experimental design (pre-registered)}
\label{sec:design}
\paragraph{Pool.} BigCodeBench (Complete split, fixed revision), restricted to tasks whose canonical solution passes the official hidden
tests in three independent runs inside the official evaluation image and whose visible doctests can be executed: 403 eligible tasks
after removing six development tasks. \draftnote{eligibility: gate1 411, gate2 408, history-stable 407, eligible 403 --- cite
bcb\_eligibility.json.}

\paragraph{Splits.} Exploration $N_e=100$ = the 59 re-eligible tasks of an earlier 60-task probe (their first-pass candidates are
\emph{reused} as $s_0$, without selection on outcome, with the original model id, prompt and cost preserved) + 41 tasks drawn with a
fixed seed; confirmation $N_c=200$ drawn with the same seed from the 344 unused tasks and never touched before the frozen policy is
applied; 103 sealed reserve tasks.

\paragraph{Models and arms.} $A$ = claude-haiku-4-5, $B$ = claude-sonnet-4-6, temperature 0.2, 4096 output tokens; six arms with one
repetition ($k=1$); a 30-task exploration subset additionally runs A-self and B-expert twice to report flip rates. A round-0 pass
measures the median cost of one B-expert call to set $b_{\mathrm{cont}}$, the per-arm continuation budget that bounds
\texttt{max\_tokens} on every attempt.

\paragraph{Pool gate.} Confirmation proceeds only if the best exploration arm success rate is $\le 80\%$. The original registered
ceiling for $B$ alone was 70\%; the probe measured $B$-solo at 75.0\% (45/60, 95\% CI [62.8, 84.2]), which was recorded as
\emph{unconfirmed} under that rule, and the operating criterion was revised to 80\% (headroom $\ge 2\times$ the 10~pp practical
threshold) \emph{before} the exploration set was assembled. Both values and the revision history are kept.

\paragraph{Primary rule.} On the 200 confirmation tasks, $d_i=y_{i,\hat\pi_{\text{post}}(i)}-y_{i,\hat a}$; exact one-sided McNemar
test on discordant pairs ($p<0.05$) \emph{and} mean $d\ge 10$~pp $\Rightarrow$ \emph{pass}, otherwise \emph{unconfirmed}. Passing
means ``evidence of positive gain and a point estimate above the practical threshold''; it does \emph{not} establish that the true
gain exceeds 10~pp. Unconfirmed is not ``no gain''. The rule, thresholds and sets are hashed into the approval record.

\paragraph{Secondary analyses.} P2-post vs.\ P2-pre paired difference (an ablation of post-execution state, each policy tuned on its
own menu), P1, per-arm success/cost/refusal, the arm distribution chosen by the policy, and cost from two viewpoints: total experiment
spend (every arm executed on every task, from the budget ledger) and \emph{deployment cost} of the policy (chosen arm + $s_0$
acquisition, charged even when the policy chooses B-solo; the fixed B-solo baseline pays only its own call).

\paragraph{Operating characteristics (synthetic).} With the identical exploration procedure applied to a synthetic generator
(100 exploration / 200 confirmation tasks, 20 seeds per condition, 60 per null), both null conditions---features independent of
outcomes, and features that predict only difficulty with an invariant best arm---passed in 0/60 cells; observed pass rates at a
calibrated Bayes-optimal gain of 15~pp were 16/20 with 6 features and fell with menu size, and at 10~pp were $\le 3/20$.
\draftnote{Table from evidence/week2/pilot\_power\_v4.md; state that these are observed pass rates under those synthetic conditions, not
power statements; gate hold rates are a property of the generator.}

\section{Results}
\label{sec:results}
\draftnote{Exploration not yet executed (approval pending). Tables are generated by paper/make\_tables.py from evidence/week2 artifacts;
nothing here is filled by hand.}
\subsection{Exploration set}
\draftnote{arm table; pool gate; frozen policy; spend from ledger.}
\subsection{Confirmation set (primary)}
\draftnote{$d$, $b$, $c$, McNemar $p$, bootstrap CI of $d$, verdict.}
\subsection{Secondary and cost}
\draftnote{post vs pre; P1; pick distribution; experiment spend vs deployment cost; measured vs conservatively attributed unknown cost.}

\section{Related work}
\label{sec:related}
\textbf{Cascades with a learned quality scorer.} FrugalGPT~\citep{chen2023frugalgpt} (Sec.~3, ``LLM cascade'') scores a generated answer
with a trained regression model and returns it if the score exceeds a threshold, otherwise queries the next API. We share the
``look at the generation, then decide'' structure, but our signal is an executed check rather than a learned scorer, the action set
includes revision by the same model and review by another, and we account for the cost of acquiring the state.
\textbf{Belief-guided delegation.} REDEREF~\citep{hosseini2026rederef} routes recursively without training, updating per-agent beliefs
from a programmatic or LLM judge, re-routing after reflection-driven critique and terminating on success or budget. Our P2-pre is
\emph{not} a re-implementation or proxy of REDEREF; the pre/post comparison is an ablation of post-execution state within our own
learner, and a direct comparison to REDEREF is future work.
\textbf{Knowledge-aware expert selection.} KABB~\citep{zhang2025kabb} selects a subset of experts with knowledge-distance priors and
Thompson sampling updated from task outcomes. \textbf{When collaboration helps.} \citet{kim2026outgrow} vary coordination structure and
capability across 260 configurations and find single-agent baseline strength the most robust predictor of whether coordination helps;
our question is the task-level, state-conditional version of theirs for one coding benchmark. \textbf{Sequential inspection and
next-hop routing.} Pandora-style routing~\citep{fisch2026pandora} treats expert inspection as costly sequential search; BiRouter~\citep{yang2026birouter}
learns a next-hop (including a finisher) from local state. \draftnote{Confirm each citation against
docs/related\_work\_week1.md: which were read in full (arXiv HTML) vs abstract only; Pandora/BiRouter entries are abstract+partial.}

\section{Limitations}
\label{sec:limits}
One benchmark, one model pair, one repetition per arm; the policy is learned from 100 tasks, and the synthetic study shows that the
learned policy captures only part of the Bayes-optimal gain, increasingly so with larger feature menus. Passing the primary rule does
not certify a $\ge 10$~pp true gain. Hidden tests of BigCodeBench are public, so contamination cannot be excluded for either model.
Visible verification has zero API cost but non-zero compute and latency, which we report as wall time, not as \$0.
\draftnote{Add: gate-hold behaviour, incomplete/unknown-cost attribution, what an unconfirmed result would and would not mean.}

\subsubsection*{Reproducibility statement}
All pre-registration documents, configuration, ID files, the budget ledger, and the exact commands are in the anonymised repository
snapshot submitted as supplementary material. \draftnote{Anonymise paths, run ids, and remove the git remote before zipping; check PDF
metadata.}

\subsubsection*{AI use statement}
\draftnote{Required section (not counted). Fill per ICLR 2027 AI Policy for Authors: generative AI tools were used to draft code,
tests and text under the authors' direction; all code was tested and all claims were checked against artifacts by the authors. The
authors must confirm this wording themselves.}

\bibliography{main}
\bibliographystyle{iclr2027_conference}

\appendix
\section{Pre-registration and approval record}
\draftnote{Reproduce PREREG\_pilot\_v4.md config block and APPROVAL.json fields (anonymised).}
\section{Synthetic operating characteristics}
\draftnote{Full table from pilot\_power\_v4.md, three quantities (a)/(b)/(c) kept separate; aggregated truth-bin rows without CIs.}
\end{document}
